\section{Other agent design patterns}
\label{sec:patterns}
\paragraph{Learning objectives.}
After this section you will be able to name the design patterns in Google Cloud's guide
\citep{google2026patterns} and sort them by one criterion: whether a model or predefined
logic decides which agent runs next. You will be able to state, for the coordinator,
hierarchical decomposition, parallel, and review-and-critique patterns, what changes in
task ownership, state boundaries, model-call overhead, latency, failure propagation, and
inspectability relative to the serial agent of \cref{sec:reflection}. And you will be able
to answer the question that closes the block: given forty alerts, a fixed per-alert budget,
and an evidence requirement, which of the five versions you would deploy, and which axis decides it.

\subsection{The problem the reading addresses}

The serial agent built in \cref{sec:fixed,sec:reactive,sec:planning}
is one point in a design space, and the space has names. The guide this section uses is
supplementary engineering reading rather than study material: it is a vendor's
architecture-center page, last updated 2026-05-28, whose purpose is to help a team choose
an orchestration shape for a workload. It is cited here for two reasons. It names the
patterns that appear in framework documentation and in job descriptions, so students will
meet them. And its own classification of the patterns rests on the question this block has
asked at every version: does a model orchestrate, or does predefined logic?

The guide opens with four requirements it says a team must settle before choosing. Does the
task decompose into predefined workflow steps, or is it open-ended and in need of a model to
orchestrate? Does the workload prioritize fast responses, or can it tolerate delay for
accuracy? What is the inference budget, and does it support several model calls per
request? Are decisions high-stakes enough to need a person in the loop? It
adds a caution the course endorses: a predictable, structured workload that a single model
call can serve does not need an agentic system at all. Those four questions are the
syllabus's axes in a vendor's words, and \cref{tab:patterns-axes} returns to them.

\subsection{Mechanism}

The guide describes twelve patterns. \Cref{tab:patterns-catalogue} lists them with the one
property this course cares about first, the orchestration decision, and the tradeoff the
guide itself states for each. Three of the multi-agent patterns (sequential, parallel,
loop) are run by what the guide calls \emph{workflow agents}, which operate on predefined
logic without consulting a model for orchestration; the coordinator, hierarchical, and
swarm patterns use a model to route or decompose. The review-and-critique and iterative
refinement patterns are, in the guide's words, implementations of the loop pattern.

\begin{table}[htbp]
  \caption{The guide's twelve patterns, sorted by who decides which agent runs next.
  Tradeoffs are the guide's own statements, paraphrased.}
  \label{tab:patterns-catalogue}
  \small
  \begin{tabularx}{\linewidth}{>{\RaggedRight}p{3.1cm}>{\RaggedRight}p{2.4cm}XX}
    \toprule
    Pattern & Orchestrated by & What it is & Tradeoff the guide states \\
    \midrule
    Single-agent & The model, within one agent & One model, one tool set, one system prompt; the model plans and picks tools & Degrades as tools and task complexity grow: latency, wrong tool choice, incomplete tasks \\
    Sequential & Predefined logic & Specialized agents in a fixed linear order; each output is the next input & Lower latency and cost than model orchestration; rigid; cannot skip unneeded steps \\
    Parallel & Predefined logic & Independent subagents run at once; a gather step synthesizes & Lower latency; higher immediate resource and token use; the gather step must reconcile conflicting results \\
    Loop & Predefined logic & A sequence of subagents repeats until an exit condition (iteration cap or a state) & Risk of an unbounded loop if the exit condition is wrong or never reached \\
    Review and critique & Predefined logic (a loop) & A generator produces; a critic evaluates against criteria and approves, rejects, or returns feedback & At least one extra model call per evaluation; latency and cost accumulate with each revision \\
    Iterative refinement & Predefined logic (a loop) & One or more agents improve a result held in session state until a quality threshold or an iteration cap & Latency and cost per cycle; exit conditions must be designed \\
    Coordinator & The model & A central agent decomposes the request and dispatches subtasks to specialists & More model calls than a single agent; higher token throughput, cost, and latency \\
    Hierarchical task decomposition & The model, at several levels & A root agent decomposes; subagents decompose further until workers can execute & Considerable architectural complexity; harder to design, debug, and maintain; many model calls \\
    Swarm & No central orchestrator; agents hand off to one another & All-to-all communication among specialists; a dispatcher starts the task; an explicit exit condition is required & The most complex and costly; risk of unproductive loops or failure to converge \\
    Reason and act (ReAct) & The model, within one agent & Thought, action, observation repeated until an answer, an iteration cap, or an error & Higher end-to-end latency than one query; a misleading tool result can propagate to the answer \\
    Human-in-the-loop & Predefined checkpoints & The agent pauses at a checkpoint and waits for a person to approve, correct, or supply input & Requires building and maintaining the external interaction system \\
    Custom logic & Code the team writes & Conditional orchestration in code, mixing the other patterns & The team designs, implements, and debugs the whole flow \\
    \bottomrule
  \end{tabularx}
\end{table}

Four of these are the ones the lecture discusses, each with the question the teaching plan
attaches to it.

\begin{description}[style=nextline]
  \item[Coordinator and specialist routing.] Who selects the worker, what context is passed
  to it, and who owns the returned result? In the guide's pattern the coordinator is a
  model call that analyzes the request and routes it; the guide contrasts this with the
  parallel pattern, whose dispatch is hardcoded.
  \item[Hierarchical task decomposition.] How are subgoals assigned, and how are partial
  results integrated? The guide presents this as the coordinator pattern applied
  recursively, for ambiguous problems that need research, planning, and synthesis.
  \item[Parallel work.] Which tasks are independent, and what does their eventual join
  require? The guide's example fans one input out to four specialists and gathers four
  outputs into one analysis.
  \item[Review and critique, and iterative refinement.] Where does feedback route control,
  and what causes the system to stop? The guide's critic can approve, reject, or return
  the output with feedback; its refinement loop stops at a quality threshold or an
  iteration cap.
\end{description}

The guide's own account of the ReAct pattern deserves one sentence, because it differs in
emphasis from \cref{sec:reactive}: it describes the observation as something the model
``saves in its memory'', and it lists an error that prevents continuation as a terminating
condition. Both are correct as descriptions of what a framework does; in this course's
vocabulary the runtime records the observation and the runtime declares the terminal
status, and \cref{sec:reactive} said why that ownership matters.

\subsection{Evidence}

The guide reports no experiment, and none is claimed here. What it offers is a comparison
that sorts the patterns by workload, reproduced in \cref{tab:patterns-comparison}. The
sorting is useful because it makes the guide's implicit criterion visible: the first group
is exactly the patterns whose orchestration is predefined, the second is exactly the
patterns whose orchestration is a model, and the third is the patterns that loop. That is
the design-time versus run-time question, and the termination question, applied to
multi-agent systems. The academic papers of \cref{sec:fixed,sec:reactive,sec:planning}
remain the evidence for what each shape costs and buys; the guide is a map of the names.

\begin{table}[htbp]
  \caption{The guide's comparison section, by workload type. The characteristics are the
  guide's, paraphrased.}
  \label{tab:patterns-comparison}
  \small
  \begin{tabularx}{\linewidth}{>{\RaggedRight}p{3.4cm}>{\RaggedRight}p{3.3cm}X}
    \toprule
    Workload & Patterns the guide lists & Characteristics the guide gives \\
    \midrule
    Deterministic & Sequential; parallel; iterative refinement & Steps known in advance; no model orchestration; a fixed sequence, or independent subtasks, or progressive improvement over cycles \\
    Dynamic orchestration & Single agent; coordinator; hierarchical decomposition; swarm & The system plans, delegates, and coordinates without a script; higher latency and cost from the orchestrating model calls \\
    Iteration & ReAct; loop; review and critique; iterative refinement & Cycles of refinement, feedback, and improvement; a distinct validation step; unpredictable latency while waiting for an exit condition \\
    Special requirements & Human-in-the-loop; custom logic & Human supervision for high-stakes or subjective tasks; branching logic beyond a linear sequence \\
    \bottomrule
  \end{tabularx}
\end{table}

\subsection{Which decisions the programmer made}

\Cref{tab:patterns-axes} takes the four selected patterns and asks what each would change
about the serial agent, axis by axis. The serial agent is the reference row: one process,
one state record, one investigation at a time, model calls bounded by the attempt limit
and the overall execution limit, failures contained within an attempt, and a trace that is
one sequence.

\begin{table}[htbp]
  \caption{What each pattern changes relative to the serial agent of \cref{sec:reflection}.}
  \label{tab:patterns-axes}
  \small
  \begin{tabularx}{\linewidth}{>{\RaggedRight}p{2.5cm}XXXX}
    \toprule
    Axis & Coordinator & Hierarchical & Parallel & Review and critique \\
    \midrule
    Task ownership & The coordinator owns the task and delegates pieces; the specialist owns only its piece & Ownership is nested; each level owns the decomposition it made & Each worker owns one independent piece; the gather step owns the whole & Unchanged: the generator owns the candidate, the critic owns acceptance \\
    State boundaries & The coordinator's context and each specialist's are separate; what crosses is what the coordinator passes & A state boundary per level; partial results cross upward & One state record per worker; the join reads all of them & Unchanged: one record; the critic's reasons are added to it \\
    Model-call overhead & One routing call per request plus the specialists' calls & Multiplied by the depth of the hierarchy & The sum of the workers' calls, run concurrently & One evaluator call per candidate, plus revisions \\
    Latency & The routing call is on the critical path & Each level adds a round trip & Reduced to the slowest worker plus the join & Each revision adds a round trip \\
    Failure propagation & A misrouted task fails in the wrong specialist; the coordinator may not notice & A wrong decomposition propagates down and its cost multiplies & Contained per worker until the join, where a missing piece blocks the whole & A critic that accepts a wrong candidate passes it on; one that rejects a right one costs a revision \\
    Inspectability & The routing decision and its reasons are in the trace, if recorded & Traces nest; which level decided what must be recoverable & Traces interleave; which worker wrote what must be recoverable & Improved: acceptance carries reasons \\
    \bottomrule
  \end{tabularx}
\end{table}

Three things the table makes visible are worth stating. The two patterns that use a model
to orchestrate (coordinator, hierarchical) move a control decision from design time to run
time. Everything the course said about the loop applies to that decision: it is adaptive, its
cost is not known in advance, and it must be recorded to be inspected. The parallel
pattern changes no control decision at all; it changes who runs when, and its whole cost
is at the join. That is Unit 2's subject, and HW2's. And the review-and-critique pattern is
not new: it is the validator and the reflection step under another name. The guide's own tradeoff
statement, that latency and cost accumulate with each revision, is the reason
\cref{sec:reflection} put an attempt limit and an overall execution limit on the loop.

\subsection{The demo}

There is no sixth version. The demonstration for this section is to take the serial graph
of \cref{sec:reflection} and say, for each selected pattern, which node would split and
which edge would appear.

\begin{itemize}
  \item \textbf{Coordinator.} The planner becomes a router. Instead of producing a plan
  for one executor, it selects among executors specialized by alert kind, and the edge from
  the planner fans out to whichever it chose. Task ownership passes to the chosen executor
  for the duration of its run and returns to the router with the candidate finding. The
  routing decision is a model call and must be traced as one.
  \item \textbf{Hierarchical decomposition.} The planner's subgoals stop being advisory
  entries in one state record and become dispatched tasks, each with its own executor and
  its own state. A new join node integrates the partial results before the validator sees a
  candidate. The graph gains a level, and so does the trace.
  \item \textbf{Parallel work.} Nothing inside the serial agent changes. Around it, forty
  alerts fan out to forty instances of the same graph, and a report node joins their
  findings. The join is where independence must be checked: two alerts that share a helper
  file are not independent in their evidence, only in their verdicts. This is HW1's report
  stage and, with bounded concurrency, HW2's experiment.
  \item \textbf{Review and critique; iterative refinement.} Already present. The validator
  is the critic, the reflection node and the reset are the refinement, and the two limits
  are the exit conditions. The one change the guide's version would make is to let the
  critic return feedback to the same generator without a reset; \cref{sec:reflection} said
  why this course resets the attempt and retains only the guidance.
\end{itemize}

\begin{figneeded}{Google Cloud design-pattern diagrams}
One SVG per pattern on the guide's page \citep{google2026patterns}, content licensed CC BY
4.0. Credit each as ``Google Cloud Architecture Center, \emph{Choose a design pattern for
your agentic AI system}, [pattern] diagram, CC BY 4.0.'' Every path below is under
\url{https://docs.cloud.google.com/static/architecture/images/} and begins with
\texttt{choose-design-pattern-agentic-ai-system-}.
\begin{center}\small
\begin{tabular}{ll}
  \toprule
  Pattern & File name suffix \\
  \midrule
  Coordinator & \texttt{coordinator.svg} \\
  Hierarchical task decomposition & \texttt{hierarchical-task.svg} \\
  Parallel & \texttt{parallel.svg} \\
  Review and critique & \texttt{review-critique.svg} \\
  Iterative refinement & \texttt{iterative-refinement.svg} \\
  ReAct & \texttt{react.svg} \\
  Loop & \texttt{loop.svg} \\
  \bottomrule
\end{tabular}
\end{center}
\end{figneeded}

The block closes on the question that has the shape of a midterm question. Given forty
alerts, a fixed per-alert budget, and a requirement that every verdict carry evidence
references that resolve: which of the five versions would you deploy, and which of the four axes,
adaptability, model-call overhead, failure propagation, inspectability, decides it? The
answer is not in these notes; the argument for it is. The chain cannot follow
evidence it did not anticipate, so it fails the evidence requirement on any alert whose
helper delegates. The loop can, at a cost that is not known before the run,
which the per-alert budget must cap. The validator is what makes the evidence
requirement checkable rather than hoped for. Whether the reflection step earns its second
attempt under a fixed budget is an empirical question the trace can answer, and it is the
question HW1 asks teams that adopt it.

Part A ends here. After the break, Part B takes up the execution machinery that
\cref{sec:reactive} only pointed at: the state schema and its reducers, transitions and
conditional routing as predicates over state, terminal states, tool signatures and the
request, dispatch, observation, state-update cycle, and execution authority.

\subsection{Exercises}

\begin{enumerate}
  \item Using \cref{tab:patterns-catalogue}, sort the twelve patterns into those where a
  model decides which agent runs next and those where predefined logic does. For each
  model-orchestrated pattern, name the trace event that would have to exist for the
  routing decision to be inspectable.
  \item The guide lists ``iterative refinement'' under both deterministic workflows and
  workflows that involve iteration (\cref{tab:patterns-comparison}). Explain, using the
  distinction between decision timing and decision ownership from \cref{sec:fixed}, why
  a loop with a fixed exit condition can belong in both.
  \item Take the parallel row of \cref{tab:patterns-axes}. Two of the forty alerts share
  one helper file, \demoref{helpers/\allowbreak separate\_request.py}. Say what the join must check
  before it treats their findings as independent, and what the trace of each instance must
  record to make that check possible.
  \item Answer the closing question for a budget that permits, on average, one attempt per
  alert. Which version do you deploy, which axis decided it, and what single number from
  a trace would change your answer?
\end{enumerate}

\subsection{Reading}

Supplementary, cited in lecture: \citet{google2026patterns}, the requirements section, the
pattern descriptions, and the comparison section. The academic papers of
\cref{sec:fixed,sec:reactive,sec:planning} and AIMA remain the
principal study material for this block.
